\documentclass[../../../include/open-logic-section]{subfiles}

\begin{document}

\olfileid{pl}{prp}{com}

\olsection{命题逻辑的完备性}

\begin{defn}
\ollabel{def:MaxCon}
公式集合 $\Gamma$ 是\emph{极大一致的}，如果它是一致的，并且对于任意一致的集合 $\Delta$，若 $\Gamma
\subseteq \Delta$，则 $\Gamma = \Delta$。
\end{defn}

\begin{lem}[真值引理]
  \ollabel{lem:truth}
设 $\Gamma$ 是极大一致的，则：
\begin{enumerate}
\item \ollabel{lem:truth-1} $!A \in \Gamma$ 当且仅当 $\lnot !A
  \notin \Gamma$；
\item \ollabel{lem:truth-2} $\Gamma \Proves !A$ 当且仅当 $!A \in \Gamma$；
\item \ollabel{lem:truth-3} $!A \lif !B \in \Gamma$ 当且仅当
  $!A \notin \Gamma$ 或 $!B \in \Gamma$。
\end{enumerate}
\end{lem}

\begin{proof}
\begin{enumerate}
\item 假设 $!A \in \Gamma$。如果也有 $\lnot !A \in \Gamma$，那么
  $\Gamma$ 是不一致的；而如果 $!A$ 和 $\lnot!A$ 都不在 $\Gamma$ 中，则由 \olref[axd]{prop:phi}，$\Gamma\cup\{!A\}$ 和
  $\Gamma \cup \{\lnot !A\}$ 中至少有一个是一致的，这意味着 $\Gamma$
  不是极大的。另一个方向的证明同理。
\item 如果 $!A \in \Gamma$，那么显然 $\Gamma \Proves !A$。另一方面，假设 $\Gamma \Proves !A$。如果 $!A \notin \Gamma$，那么由 \olref{lem:truth-1}，有 $\lnot !A \in \Gamma$。于是有 $\Gamma \Proves
  !A$ 且 $\Gamma \Proves \lnot !A$，得出矛盾。
\item 练习。
\end{enumerate}
\end{proof}

\olref{lem:truth}\olref{lem:truth-2} 的从左到右方向表明 $\Gamma$ 是\emph{演绎封闭的}，即对所有 $!A$，若 $\Gamma \Proves !A$，则 $!A
\in \Gamma$。

\begin{prob}
  完成 \olref{lem:truth} 的证明。
\end{prob}

\begin{thm}[完备性定理]
\ollabel{thm:completeness}
如果 $\Gamma$ 是一致的，那么 $\Gamma$ 是可满足的。
\end{thm}

\begin{proof}
令 $!A_0$, $!A_1$,~\dots 为该语言所有公式的一个穷尽列表。我们递归地定义一个递增的公式集合序列 $\Gamma_0$, $\Gamma_1$,~\dots 如下：
% Part: propositional-logic
% Chapter: propositional-logic
% Section: completeness
然后定义：
\begin{align*}
\Gamma_0 & = \Gamma\\
\Gamma_{n+1} &  =
\begin{cases}
  \Gamma_n \cup \{!A_n\} & \text{如果 $\Gamma_n \cup \{!A_n\}$ 是
    一致的；}\\
  \Gamma_n \cup \{\lnot !A_n\} & \text{否则}.
\end{cases}
\end{align*}
\[
\Gamma^* = \bigcup_{0\le n}\Gamma_n.
\]
接下来的证明通过依次确立以下事实来完成：
\begin{enumerate}
\item 对每个 $n$，集合 $\Gamma_n$ 是一致的（通过对 $n$ 归纳，使用 \olref[axd]{prop:phi}）；
\item $\Gamma^*$ 是一致的；
\item $\Gamma^*$ 是极大的。
\end{enumerate}
然后定义赋值 $v$，规定 $v(p_i) = \True$ 当且仅当 $p_i \in \Gamma^*$。接着通过对 $!A$ 进行归纳，可以证明对于更复杂的语句，公式属于 $\Gamma^*$ 当且仅当它在赋值 $v$ 下为真：$\overline{v}(!A) = \True$ 当且仅当 $!A \in \Gamma^*$。特别地，$v$ 满足 $\Gamma^*$，且由于 $\Gamma \subseteq \Gamma^*$，$\Gamma$ 也是可满足的，如所欲证。
\end{proof}

\begin{cor}
如果 $\Gamma \Entails !A$，那么 $\Gamma \Proves !A$。
\end{cor}

\begin{proof}
如果 $\Gamma\Proves/ !A$，那么根据 \olref[axd]{prop:prov-incons}，$\Gamma \cup \{\lnot !A\}$ 是一致的。因此由完备性定理，$\Gamma \cup \{\lnot !A\}$ 是可满足的，从而 $\Gamma \Entails/ !A$。
\end{proof}

\begin{prop}[紧致性定理]
\ollabel{prop:compactness}
$\Gamma$ 是可满足的，当且仅当 $\Gamma$ 的每个\emph{有限}子集 $\Gamma_0$ 是可满足的。
\end{prop}

\begin{proof}
$\Gamma$ 不可满足当且仅当它不一致，当且仅当 $\Gamma$ 的某个有限子集 $\Gamma_0$ 不一致，当且仅当 $\Gamma$ 的某个有限子集 $\Gamma_0$ 不可满足。
\end{proof}

\begin{cor}
$\Gamma \Entails !A$ 当且仅当存在 $\Gamma$ 的某个有限子集 $\Gamma_0$，使得 $\Gamma_0 \Entails !A$。
\end{cor}

\end{document}